\chapter*{Mở đầu}
\addcontentsline{toc}{chapter}{Mở đầu}
\markboth{Mở đầu}{Mở đầu}

\section*{Đặt vấn đề}
\addcontentsline{toc}{section}{Đặt vấn đề}
\markboth{Đặt vấn đề}{Đặt vấn đề}

Trong bối cảnh chuyển đổi số quốc gia đang diễn ra mạnh mẽ,
bảo đảm an toàn thông tin (ATTT) cho các hệ thống công nghệ
thông tin (CNTT), đặc biệt là tại các doanh nghiệp, cơ quan
nhà nước và hạ tầng trọng yếu của quốc gia, là yêu cầu quan
trọng và cấp bách. Các cuộc tấn công mạng ngày càng tinh vi,
có chủ đích (APT - Advanced Persistent Threat) và thường sử
dụng các mẫu phần mềm độc hại (malware) được tùy chỉnh riêng
biệt, chưa từng xuất hiện trước đó trên thế giới. Điều này
khiến các hệ thống phòng thủ truyền thống, như antivirus
dựa trên chữ ký (signature-based), trở nên kém hiệu quả
trước các mối đe dọa mới. Theo thống kê mới nhất, các hệ
thống an ninh mạng phát hiện trung bình hơn 467.000 mẫu phần
mềm độc hại mỗi ngày trên toàn cầu trong năm 2024, tăng mạnh
so với các năm trước. \cite{kaspersky2024} Tại Việt Nam, 46,15\% cơ quan và doanh
nghiệp bị tấn công mạng ít nhất một lần, với tổng hơn 659.000
vụ tấn công được ghi nhận, trong đó 26,14\% là các cuộc tấn
công có chủ đích (APT) sử dụng mã độc nằm vùng, và khoảng
14,5\% bị tấn công bằng ransomware mã hoá dữ liệu. Trên thực
tế, gần 75\% chiến dịch APT toàn cầu có sử dụng malware, và
có đến 80\% trong số đó tồn tại trong hệ thống hơn 6 tháng
trước khi bị phát hiện, \cite{safegate2024} cho thấy rõ ràng các phương pháp
phòng thủ truyền thống dựa trên chữ ký đang ngày càng kém
hiệu quả trước các mối đe dọa tinh vi, mới xuất hiện.

Do đó, việc nghiên cứu và triển khai một mô-đun phát hiện
mã độc định dạng Windows Portable Executable (PE) dựa trên
học máy (machine learning) là yêu cầu cấp thiết, đồng thời
mang ý nghĩa khoa học và thực tiễn rõ rệt. Khác với các
phương pháp truyền thống dựa trên chữ ký, các mô hình học máy
có khả năng tự động học đặc trưng từ dữ liệu, cho phép phân
tích sâu các đặc điểm tĩnh (cấu trúc file, header PE, bảng
import/export, entropy, chuỗi ký tự\dots) cũng như hành vi của
mã độc. Nhờ đó, hệ thống có thể phát hiện hiệu quả các mẫu
mã độc mới, chưa từng xuất hiện trước đây (zero-day malware),
vốn là điểm yếu lớn của các giải pháp antivirus truyền thống.

Việc áp dụng học máy trong bài toán phát hiện mã độc không
chỉ giúp nâng cao tỷ lệ phát hiện (detection rate) mà còn
góp phần giảm thiểu tỷ lệ cảnh báo sai (false positive),
từ đó nâng cao độ tin cậy của các hệ thống phòng thủ an
toàn thông tin. Đặc biệt, trong bối cảnh các chiến dịch
tấn công có chủ đích (APT) ngày càng phổ biến và thường
sử dụng mã độc được tùy chỉnh riêng cho từng mục tiêu,
khả năng phát hiện dựa trên hành vi và đặc trưng tổng quát
của các mô hình học máy đóng vai trò quan trọng trong việc tăng
cường chiều sâu phòng thủ (defense in depth) cho các hệ
thống CNTT.

Bên cạnh ý nghĩa về mặt kỹ thuật, đề tài còn có giá trị
nghiên cứu và đào tạo rõ rệt. Việc xây dựng mô-đun phát
hiện mã độc PE dựa trên học máy giúp người nghiên cứu
tiếp cận một cách hệ thống các kiến thức liên ngành,
bao gồm: kiến trúc hệ điều hành Windows, định dạng file PE,
kỹ thuật phân tích mã độc, xử lý dữ liệu, cũng như các thuật
toán học máy hiện đại. Qua đó, đề tài góp phần củng cố nền
tảng lý thuyết và kỹ năng thực hành trong lĩnh vực an toàn
thông tin và trí tuệ nhân tạo - vốn là những lĩnh vực then
chốt trong khoa học máy tính hiện nay.

Ngoài ra, đề tài còn mang ý nghĩa quan trọng trong việc:

\begin{itemize}
    \item Góp phần nâng cao năng lực nghiên cứu và tự chủ
    công nghệ trong lĩnh vực an ninh mạng, thông qua việc
    chủ động xây dựng và đánh giá các mô hình phát hiện
    mã độc, thay vì hoàn toàn phụ thuộc vào các giải pháp
    đóng và công nghệ nước ngoài.
    \item Đáp ứng yêu cầu bảo đảm an toàn thông tin trong
    quá trình chuyển đổi số tại Việt Nam, phù hợp với các
    định hướng và chiến lược phát triển an toàn, an ninh mạng
    do Chính phủ đề ra, \cite{lsvn2026} đặc biệt trong bối
    cảnh gia tăng các mối đe dọa mạng nhắm vào hệ thống
    CNTT trong nước.
\end{itemize}

Với những lý do trên, đề tài "Ứng dụng học máy trong phát hiện
mã độc PE trên Windows" là một hướng nghiên cứu có giá trị
khoa học, tính thời sự và khả năng ứng dụng thực tiễn cao,
đồng thời phù hợp với mục tiêu đào tạo và nghiên cứu của
các chuyên ngành liên quan đến an ninh, an toàn thông tin
trong các cơ sở giáo dục đại học hiện nay.


















\section*{Mô tả bài toán}
\addcontentsline{toc}{section}{Mô tả bài toán}
\markboth{Mô tả bài toán}{Mô tả bài toán}

Khóa luận tập trung giải quyết hai bài toán có liên hệ mật thiết với
nhau: \textit{bài toán nghiên cứu} phương pháp học máy để phân loại file
PE (phân loại nhị phân, với nhãn phân loại là lành tính hoặc mã độc),
và \textit{bài toán công nghệ} về triển khai dịch vụ phát hiện mã độc trên
máy endpoint (Windows).

Cụ thể, bài toán nghiên cứu mà khóa luận giải quyết:

\begin{itemize}
    \item \textbf{Đầu vào:} Một file thực thi định dạng PE trên Windows,
        cùng với nhãn phân loại tương ứng (trong giai đoạn huấn luyện).
    \item \textbf{Đầu ra:} Nhãn phân loại (lành tính - benign hoặc
        mã độc - malware).
    \item \textbf{Phạm vi:} Phân tích tĩnh (static analysis), không
        thực thi file trong môi trường sandbox. Dữ liệu thực nghiệm sử
        dụng tập EMBER2024.
    \item \textbf{Giả định:} File PE hợp lệ về mặt cấu trúc và có thể
        đọc được bằng các công cụ phân tích tĩnh thông thường.
    \item \textbf{Phương pháp:} Xây dựng và so sánh các pipeline học
        máy kết hợp đặc trưng từ EMBER2024 với đặc trưng năng lực mã
        độc được trích xuất bằng công cụ Capa. Áp dụng kỹ thuật giảm
        chiều TruncatedSVD để nén đặc trưng Capa, và SHAP để giải thích
        kết quả dự đoán của mô hình.
\end{itemize}

Bài toán công nghệ mà khóa luận giải quyết:

\begin{itemize}
    \item \textbf{Đầu vào:} Một file PE cần kiểm tra, được cung cấp
        thông qua giao diện người dùng (tải lên thủ công) hoặc được
        phát hiện tự động bởi driver khi có tiến trình thực thi trên
        máy endpoint.
    \item \textbf{Đầu ra:} Kết quả phân loại (lành tính hoặc mã độc),
        danh sách các năng lực/hành vi đặc trưng được trích xuất, giải
        thích trực quan quyết định của mô hình với kỹ thuật xAI, và
        lịch sử các lần quét được lưu trữ để tra cứu.
    \item \textbf{Phạm vi:} Xây dựng hệ thống gồm nhiều thành phần
        chạy trên máy endpoint Windows, bao gồm kernel driver (KMDF),
        Windows service, module trích xuất đặc trưng (C++), giao diện
        người dùng trên máy endpoint, và dịch vụ phân tích chuyên sâu
        online (deep-analysis). Phần triển khai trong môi trường thực tế
        (production) được tiếp cận ở mức mô phỏng và phân tích kiến
        trúc; đề tài không đặt mục tiêu triển khai hoàn chỉnh trên hệ
        thống quy mô lớn.
    \item \textbf{Giả định:} Môi trường thực thi là máy tính cá nhân
        hoặc máy trạm chạy Windows, có đủ quyền để cài đặt driver và
        Windows service. Riêng dịch vụ phân tích file PE chuyên sâu
        online được triển khai dưới dạng ứng dụng web trên hệ điều
        hành Linux.
    \item \textbf{Phương pháp:} Tích hợp pipeline học máy từ bài toán
        nghiên cứu vào kiến trúc dịch vụ gồm hai chế độ - antivirus
        offline chạy nền trên endpoint và dịch vụ phân tích sâu online
        - kết hợp các kỹ thuật lập trình hệ thống Windows như lập
        trình driver (KMDF), Windows service và Windows API.
\end{itemize}








\section*{Mục tiêu và phạm vi nghiên cứu của khóa luận}
\addcontentsline{toc}{section}{Mục tiêu và phạm vi nghiên cứu của khóa luận}
\markboth{Mục tiêu và phạm vi nghiên cứu của khóa luận}{Mục tiêu và phạm vi nghiên cứu của khóa luận}

Khóa luận đặt ra ba mục tiêu chính như sau:

\begin{enumerate}
    \item \textbf{Nghiên cứu kỹ thuật trích xuất năng lực mã độc (capability extraction):}
    Tìm hiểu và làm chủ các phương pháp phân tích tĩnh file PE nhằm trích xuất
    các năng lực và hành vi đặc trưng của mã độc. Đây là nền tảng để xây dựng
    bộ đặc trưng (feature set) có ý nghĩa ngữ nghĩa cao hơn so với các đặc
    trưng cấu trúc thông thường, từ đó nâng cao hiệu quả phân loại của mô
    hình học máy.
    \item \textbf{Đánh giá thực nghiệm trên tập dữ liệu thực tế:}
    Huấn luyện và đánh giá các mô hình học máy trên tập dữ liệu EMBER2024 - một
    trong những tập dữ liệu benchmark chuẩn được công nhận, có quy mô lớn phù hợp
    cho bài toán nghiên cứu, xây dựng mô hình phát hiện mã độc PE.
    Thực nghiệm nhằm làm rõ bộ đặc trưng nào mang lại hiệu quả phân
    loại tốt nhất, cũng như đóng góp của kỹ thuật capability extraction
    vào hiệu quả tổng thể của mô hình học máy.
    \item \textbf{Thiết kế và mô phỏng kiến trúc agent bảo mật trên endpoint:}
    Nghiên cứu và thiết kế kiến trúc triển khai mô hình phát hiện mã độc dưới
    dạng một agent chạy trực tiếp trên máy endpoint Windows. Phần này tiếp cận
    ở mức mô phỏng và phân tích kiến trúc, tập trung vào việc tìm hiểu và áp
    dụng các kỹ thuật lập trình hệ thống Windows bao gồm lập trình driver (KMDF),
    lập trình Windows service và sử dụng Windows API. Đề tài không đặt mục tiêu
    triển khai hoàn chỉnh trên hệ thống quy mô lớn.
\end{enumerate}

Do đó, phạm vi nghiên cứu của khóa luận có thể được tóm tắt như sau:

\begin{itemize}
    \item \textbf{Đối tượng phân tích:} File thực thi định dạng PE (.exe, .dll)
    trên hệ điều hành Windows.
    \item \textbf{Phương pháp phân tích:} Phân tích tĩnh (static analysis);
    không bao gồm phân tích động (dynamic analysis) hay phân tích trong sandbox.
    \item \textbf{Dữ liệu:} Tập dữ liệu EMBER2024 dùng cho huấn luyện và đánh giá mô hình.
    \item \textbf{Mô hình học máy:} Tập trung vào các mô hình học máy truyền thống
    (gradient boosting).
    \item \textbf{Triển khai:} Ở mức thiết kế kiến trúc và mô phỏng thành phần
    (Proof-of-Concept), phù hợp mục tiêu nghiên cứu, áp dụng các kỹ thuật lập
    trình hệ thống trên hệ điều hành Windows.
\end{itemize}







\section*{Cấu trúc khóa luận}
\addcontentsline{toc}{section}{Cấu trúc khóa luận}
\markboth{Cấu trúc khóa luận}{Cấu trúc khóa luận}

Phần còn lại của khóa luận được tổ chức thành các chương như sau:

\textbf{Chương 1 - Cơ sở lý thuyết và công nghệ sử dụng:}
Trình bày các kiến thức nền tảng cần thiết để hiểu các phương pháp
được đề xuất ở các chương sau, bao gồm: tổng quan về mã độc và định
dạng file PE, các phương pháp phát hiện mã độc, kỹ thuật trích xuất
năng lực mã độc (capability extraction) bằng công cụ Capa, các mô hình
học máy và kỹ thuật giảm chiều dữ liệu được sử dụng, phương pháp giải
thích mô hình (xAI) với SHAP, các công trình nghiên cứu liên quan, và
các công nghệ, thư viện áp dụng trong đề tài.

\textbf{Chương 2 - Thiết kế và đánh giá các pipeline học máy trong
phát hiện mã độc:}
Trình bày kiến trúc tổng thể của giải pháp, thiết kế chi tiết ba
pipeline học máy được đề xuất dựa trên các tổ hợp đặc trưng khác nhau
(EMBER2024 thuần, EMBER2024 kết hợp Capa one-hot, EMBER2024 kết hợp
Capa TruncatedSVD), và kết quả so sánh thực nghiệm giữa các pipeline
trên tập dữ liệu EMBER2024. Từ đó xác định pipeline phù hợp nhất để
đưa vào triển khai ở Chương 3.

\textbf{Chương 3 - Phát triển dịch vụ phát hiện mã độc trên máy endpoint và dịch vụ phân tích mã độc online:}
Trình bày phân tích yêu cầu và thiết kế chi tiết hệ thống agent chạy
trên máy endpoint Windows, bao gồm các thành phần: kernel driver
(KMDF), Windows service, module trích xuất đặc trưng, server/dịch vụ
phân tích chuyên sâu online tích hợp Capa và SHAP, và giao diện người
dùng trên máy endpoint. Chương này cũng trình bày kết quả cài đặt,
các kịch bản kiểm thử và đánh giá hệ thống.

\textbf{Kết luận và hướng phát triển:}
Tổng kết các kết quả đạt được của cả hai bài toán nghiên cứu và công
nghệ, nêu rõ những hạn chế còn tồn tại, và đề xuất hướng phát triển
trong tương lai.
